% PASS20260920 physical twin-A4 index-5 E8 gluing theorem.
\subsection*{The physical twin-$A_4$ factors and the index-five $E_8$ gluing}

The recorded flagship bases contain two distinct type-$A_4$ root systems in the
same physical $E_8$: the Wilson-line $SU(5)$ root system and the
hypercharge-orthogonal $A_4$ organizing the four extra Abelian directions.
Their Gram matrices are both $C_{A_4}$ and their cross Gram matrix vanishes.
The combined rank is eight, hence
\[
\det(C_{A_4}\oplus C_{A_4})=25
\]
and, because $E_8$ is unimodular,
\[
\boxed{[E_8:A_4\oplus A_4]=5.}
\]
Thus the full-rank overlattice quotient is cyclic of order five.

The executable verifier enumerates all $240$ $E_8$ roots and computes their
orthogonal projection norms and discriminant classes $(c_C,c_G)$ in the two
$A_4$ factors.  The census is
\[
20+20+50+50+50+50=240,
\]
with mixed class pairs
\[
(1,3),\quad(2,1),\quad(3,4),\quad(4,2),
\]
so the gluing is the graph
\[
\boxed{c_G=3c_C\pmod5.}
\]
Since the nonzero $A_4$ weight classes of norms $4/5$ and $6/5$ are the
$5,\bar5$ and $10,\overline{10}$ Weyl orbits, respectively, the complete
adjoint branching recovered from the physical roots is
\[
\boxed{
248=(24,1)+(1,24)+(10,5)+(\overline{10},\overline5)
 +(\overline5,10)+(5,\overline{10}) }.
\]

\paragraph{Scope firewall.}
The complementary $A_4$ is used by the flagship as an organizer of Abelian
charges.  The lattice theorem does not assert an unbroken second $SU(5)$ gauge
factor, nor does it solve the D/F-flat vacuum equations.  It explains why the
two physical $S_5$ structures are parallel while remaining orthogonal.

The executable witness is
\texttt{analysis/w33\_physical\_twin\_a4\_e8\_index5\_gluing.py}.
